\chapter{Text, Eigenschaften und Prüfung}
\label{kap:editor-text}

\section{Der Texteditor}

Rechts neben dem Plan steht die Eingabedatei. Der Texteditor ist auf das
Format zugeschnitten:

\begin{itemize}
  \item \textbf{Hervorhebung}: Schlüssel, Zeichenketten, Zahlen und
    Kommentare in eigenen Farben.
  \item \textbf{Zeilennummern} am Rand; Zeilen mit Fehlern oder Warnungen
    sind markiert, die Meldung steht im Tooltip.
  \item \textbf{Einrückung}: Die Eingabetaste übernimmt die Einrückung der
    Zeile, nach \texttt{\{} oder \texttt{[} eine Stufe mehr; \taste{Tab} und \tasten{Umschalt}{Tab} rücken markierte Zeilen ein
    und aus.
  \item \tasten{Strg}{/} kommentiert die markierten Zeilen mit \code{//}
    aus oder wieder ein.
  \item Eine schließende Klammer am Zeilenanfang rückt von selbst aus;
    die Rückschritttaste im Einzug löscht eine ganze Einrückstufe.
\end{itemize}

Steht der Cursor in einem Bauteil, wählt der Plan dieses Bauteil aus. Steht
er in einem Bauteil einer Teilschaltung, wechselt der Plan in deren Plan.

\section{Prüfung beim Tippen}

Kurz nach jeder Änderung liest der Editor die Datei mit dem Parser des
Simulators – genau so, wie \esim\ sie lesen würde. Er baut dabei die ganze
Schaltung auf, ohne zu rechnen. Gefunden werden unter anderem:

\begin{itemize}
  \item Syntaxfehler im JSON (fehlende Kommas, Klammern, Anführungszeichen),
  \item unbekannte Schlüssel (»unknown parameter "valeu" ignored«),
  \item Bauteile mit falscher Zahl von Knoten, fehlenden Werten oder
    unbekanntem Modell,
  \item Fehler in Ausdrücken und unbekannte Namen darin,
  \item Gleichungssysteme mit falscher Zahl von Gleichungen,
  \item Widersprüche zwischen Plan und Text.
\end{itemize}

Fehler und Warnungen stehen an ihrer Zeile und gesammelt in
\menue{Probleme}; ein Doppelklick dort springt zur Zeile. Die Statuszeile
nennt die Anzahl oder, wenn alles stimmt, die Größe der Schaltung (»22
Bauteile, 18 Unbekannte«).

Nicht alles lässt sich ohne Rechnen erkennen: Ob ein Arbeitspunkt
konvergiert oder ein Knoten keinen Gleichstrompfad nach Masse hat, zeigt
erst die Simulation im \menue{Protokoll}.

\section{Eigenschaften}

Das Feld \menue{Eigenschaften} zeigt das ausgewählte Bauteil als Formular:

\begin{itemize}
  \item \textbf{Titel}: Art des Bauteils (»NPN-Transistor«,
    »Verhaltensquelle«).
  \item \textbf{Name}: Umbenennen ändert den Namen im Text. Ein Name, den
    schon ein anderes Bauteil trägt, wird abgelehnt.
  \item \textbf{Werte}: ein Feld je Schlüssel des Bauteils, dazu die
    üblichen Schlüssel seiner Art, auch wenn sie noch nicht gesetzt sind
    (grau: »nicht gesetzt«). Die Eingabetaste oder das Verlassen des Felds
    schreibt den Wert in den Text. Was gültiges JSON ist (Zahlen,
    \code{true}, Listen, ganze Objekte), wird so geschrieben, alles andere
    (\code{4.7k}, Ausdrücke) als Zeichenkette.
  \item \textbf{Neu}: Name und Wert eines weiteren Schlüssels, mit \taste{+}
    oder der Eingabetaste übernommen, z.\,B. \code{ac} = \code{1} an einer
    Quelle oder \code{ic} = \code{5} an einem Kondensator.
  \item \textbf{Anschlüsse}: die Knoten des Bauteils, zur Information.
  \item \menue{Im Text zeigen} springt zur Stelle im Text.
\end{itemize}

\paragraph{Das Feld »model«.}
Bei Dioden, Transistoren und Teilschaltungen schlägt das Feld beim Tippen
passende Modelle vor: die der Datei und die der Bibliothek, gefiltert nach
der Art des Bauteils. Darunter steht, woher das gewählte Modell kommt (»in
der Datei«, »aus der Bibliothek, wird übernommen«, »nicht definiert«). Ein
Modell aus der Bibliothek wird beim Übernehmen in \code{"models"} der
Datei kopiert. Die Schaltfläche \menue{\dots} öffnet die Modellbibliothek.

Verschachtelte Werte wie die Kurvenform \code{"wave"} einer Quelle zeigt das
Formular als Text; sie werden am besten im Texteditor bearbeitet
(\menue{Im Text zeigen}).

\section{Protokoll}

Das \menue{Protokoll} zeigt alles, was \esim\ während der Simulation
ausgibt: die Größe der Schaltung, je Analyse eine Zeile mit Schritten,
Iterationen und Ausgabedatei, den Arbeitspunkt, Messgrößen, Statistik von
Monte-Carlo-Läufen und am Ende die Rechenzeit. Warnungen und Fehler sind
farbig hervorgehoben. Scheitert eine Simulation, steht hier die Ursache.
